\section{Related Work}

\subsection{LLM in mutation testing (generation and filtering)}

The most active branch of research today uses LLMs to generate mutants. Wang et al. performed the first comprehensive study with many models (GPT-4o-mini, GPT-3.5-Turbo, DeepSeek-V3, CodeLlama-13B, StarChat-B...) on 851 real bugs in Java, including 605 bugs from 12 Defects4J 2.0 projects and 246 ConDefects bugs \cite{wang2025comprehensive}. Key result: mutants generated by LLMs detect 76.47\% of real bugs compared to 44.15\% for rule-based techniques (Major, LEAM, uBERT) - meaning an improvement of 32.99 percentage points - but the trade-off is that the non-compilable ratio was higher by 32.73 points, the duplicate ratio higher by 9.38 points, and the equivalent ratio higher by 3.64 points. SMART addressed these disadvantages using RAG on vectorized real-bug repositories, focused code chunking, and supervised fine-tuning: validity increased from 42.89\% to 65.6\%, the non-duplicate rate from 87.38\% to 95.63\%, and real-bug detection reached 92.62\% compared to 57.86\% for LLMut \cite{wang2026smart}. LLMorpheus takes a lighter approach using prompt templates \cite{tian2025leamplusplus}.

On the equivalent mutant filtering side - a well-known undecidable problem in mutation testing - Tian et al. evaluated a system of LLMs \cite{tian2024equivalent}, and GEM-LLM combined inter-procedural slicing, invariant inference via LLM, and SMT verification on 10 Defects4J projects (PIT mutants): classifying 25-30\% of surviving mutants that had been overlooked as equivalent, with 98\% precision, reaching up to 34.2\% for relational operators and 31.8\% for conditional-boundary operators, broken down by operator \cite{tanhaei2025gemllm}.

LLMs are also used to generate mutation-guided tests: MuTAP uses mutants as feedback for the program \cite{dakhel2024mutap}; MUTGEN (Llama-3.3-70B + PIT feedback) surpasses EvoSuite in mutation score, averaging 88-91\% on HumanEval-Java and LeetCode-Java \cite{wang2026mutgen}.

Common limitations of this group: in all the studies above, LLMs play a role in generating, filtering, or generating tests, but none of the papers use an LLM on a pre-existing list of mutants and ask: "with a fixed budget, which mutants should be run?"

\subsection{Predictive mutation testing and selective mutation}

PMT avoids execution costs by predicting the result: Seshat uses natural language (method name, test name) to predict the kill matrix faster than running the actual tests \cite{kim2022seshat}; MutationBERT uses CodeBERT with the full context of the mutant-test pair, significantly improving same-project and cross-project prediction \cite{jain2023mutationbert}; SODA (CodeT5 + contrastive learning) currently leads on six Defects4J v2.0.0 projects, with kill-F1 improvements of 5.32-114.92\%, survive-F1 improvements of 0.04-25.54\%, and accuracy improvements of 4.25-60.43\% compared to baselines, and compared to MutationBERT in the cross-project scenario, it achieves +11.09\% kill-F1 \cite{zhao2024soda}. On the selective mutation side, Cerebro selects subsuming mutants via static analysis \cite{garg2023cerebro}; LEAM and LEAM++ directly learn to build quality mutants instead of selecting from a pool \cite{ma2021mudelta}; AL-PMT shows that, with active learning, observing only 10\% of the kill status of the pool achieves 98\% of maximum performance in >80\% of projects - independent evidence that a 10\% budget is a meaningful working point in practice \cite{borsi2025alpmt}. Delgado-Pérez et al. summarized 8 methodological threats affecting PMT approaches as a whole - from data leakage to evaluation unit selection - and experimentally demonstrated their impact \cite{delgadoperez2026pitfalls}; our design directly incorporates these warnings (pre-registration, pairing by project/version/class, reporting both pooled units).

Common limitations of this group: these approaches need to be trained or fine-tuned per project/dataset in order to predict the result or build new mutants; they do not answer the question, "Is a training-free, zero-shot LLM subset better than random?" - a question that is meaningful for teams that lack training facilities.

\subsection{Selection compared with a random baseline}
Random selection has long been a compulsory baseline for the mutant selection problem: Zhang et al. systematically compared data evaluating mutant selection \cite{pitts2022random}, while Pitts reconsidered random selection under the influence of equivalent mutants \cite{pitts2022random}. Sentinel approached the problem via hyper-heuristics to generate reduction strategies \cite{guizzo2022sentinel}. Most directly relevant to us, Jimenez et al. compared naturalness-based ranking (n-gram) against random selection on five Defects4J projects (>100k mutants, 230 real faults) using the same statistical tools we inherited - the Wilcoxon

Signed-rank test and effect size - and obtained negative results: $A_{12} \approx 0.57-0.58$ (small), meaning naturalness is not significantly better than random selection in a way that is meaningful in practice \cite{jimenez2018naturalness}. This number is the basis for our setting the effect size threshold $r\_rb \ge 0.30$ (medium) in the proposal: an improvement at the level Jimenez observed would not be considered a success.
\subsection{Positioning}
Unlike prior work, this paper is the first to evaluate GPT-4o-mini as a
mutant-selection decision-maker against same-budget random fixed-seed selection on Defects4J,
inheriting the ranking-versus-random evaluation design of Jimenez et al.
\cite{jimenez2018naturalness} but replacing n-gram naturalness with a modern instruction-following
LLM, and additionally analyzing how the unit of aggregation (group-weighted versus
mutant-weighted) affects the statistical conclusion.